% Campionario tecnico ed editoriale: gli ambienti qui mostrati sono facoltativi.
% Pagina di prova del preambolo: aggiungi un campione delle scritture e dei simboli
% del corso, compila, guarda il risultato, poi cancella il file e la sua riga \input.
\chapter{Prova tecnica}\label{cap:prova}

\section{Scritture}
Italiano con accenti: è, à, ì, ò, ù, perché, città. Parola inglese: \en{translation}\aut{Autore, Nome}\idx{concetto}.
Greco: \gr{παρθένος}, \gr{Ἰδοὺ ἡ παρθένος ἐν γαστρὶ ἕξει}. Rimando: capitolo~\ref{cap:prova}.

\begin{cit}[english]
A quotation in the language of the source.
\fontecit{(Fonte, anno)}
\end{cit}

\afronte[english]{The source text stands on the left, in its language, across several lines of the column.}{La traduzione sta a destra, su più righe della colonna.}

\section{Riquadri}
\begin{correzione}Prova di correzione.\end{correzione}
\begin{dubbio}Prova di dubbio.\end{dubbio}
\begin{dallibro}Prova dal libro (p.~1).\end{dallibro}
\soloappunti
Testo. \integrazione{prova di integrazione.}

\section{Tabelle, formule, codice, diagrammi}
\begin{center}\small
\begin{tabularx}{\linewidth}{@{}lYY@{}}
\toprule
& \textbf{Prima colonna} & \textbf{Seconda colonna} \\
\midrule
Riga & un testo abbastanza lungo da andare a capo dentro la cella della tabella & un altro testo abbastanza lungo da andare a capo \\
\bottomrule
\end{tabularx}
\end{center}

Formula in linea $E = mc^2$, $\alpha \to \beta$, e in evidenza:
\begin{align}
\sum_{i=1}^{n} i &= \frac{n(n+1)}{2} \\
\int_0^1 x^2 \, dx &= \frac{1}{3}
\end{align}

\begin{lstlisting}[numbers=left]
int main(void) { /* commento con è */
    return 0;
}
\end{lstlisting}

\begin{center}
\begin{tikzpicture}[node distance=2.4cm,every node/.style={draw,rounded corners,font=\small,inner sep=5pt}]
\node (a) {Sorgente};
\node[right=of a] (b) {Traduzione};
\draw[-{Latex}] (a) -- (b);
\end{tikzpicture}
\end{center}

\begin{insintesi}
Prova del riquadro di fine capitolo.
\end{insintesi}
